\section{FAQ}


\begin{tcolorbox}[colframe=blue, colback=white]
	Reminder - Reading the card explains the card -- except when you have to look in here.
\end{tcolorbox}

\subsection*{General}
\begin{enumerate}
	\item If I disengage all enemies in front of me, but there are still actions on the card, can I use the remaining actions at range?\\
	\hspace*{1cm}- Yes.
	\item When I end my turn and flip the leftmost enemy sideways, does that enemy behave differently?\\
	\hspace*{1cm}- No, being sideways is just an indicator.
	\item When multiple passives trigger at the same time, what determines the order?\\
	\hspace*{1cm}- The player that owns the passives decides. You may order them as you please.
	\item Can I heal above my starting life?\\
	\hspace*{1cm}- Yes.
	\item When I roll bad enough to deal negative damage, does this heal the enemy?\\
	\hspace*{1cm}- No, the minimum damage is 0.
	\item Can I choose to not use a passive?\\
	\hspace*{1cm}- No, all effects are forced.
\end{enumerate}


\subsection*{Ancestries}

\subsubsection*{Fae}
\begin{enumerate}
	\item How does \symRanged~Ranged \symFlurry~Flurry work?\\
	\hspace*{1cm}- You can give the flurry tokens to someone else. Spending a flurry token does not benefit from your passive.
\end{enumerate}

\subsubsection*{Orc}
\begin{enumerate}
	\item Can I empower the druid 1 \symHeal~Heal with my powerful?\\
	\hspace*{1cm}- Yes, because you may order your passives.
\end{enumerate}

\subsection*{Classes}

\subsubsection*{Sangromancer}
\begin{enumerate}
	\item Can I send myself to the backline by paying life until I hit 0?\\
	\hspace*{1cm}- Yes.
\end{enumerate}

\subsubsection*{Pyromancer}
\begin{enumerate}
	\item Does "Mortal Flame" only count the initial damage or also modifiers when tripling damage?\\
	\hspace*{1cm}- All modifiers apply, such as dice, powerful, vulnerable or passives.
\end{enumerate}


\subsection*{Enemies}
\subsubsection*{Bandit Leader}
\begin{enumerate}
	\item How does stun work with his extra action?\\
	\hspace*{1cm}- Stunning the bandit leader will remove the bonus action from his adjacent bandits. Stunning an adjacent bandit will remove only one of that bandit's actions.
\end{enumerate}

\subsubsection*{Orc}
\begin{enumerate}
	\item How does powerful work on enemies?\\
	\hspace*{1cm}- It still empowers their \symAttack~Attacks, \symHeal~Heals, and \symPoison~Poisons.
\end{enumerate}
